\documentclass[11pt]{article}
\usepackage[T1]{fontenc}
\usepackage{lmodern,amsmath,amssymb,amsthm,booktabs,array}
\usepackage[margin=1in]{geometry}
\usepackage[hidelinks]{hyperref}
\hypersetup{
  pdftitle={Exact counts and sharp growth for OEIS A222759, first comment (12 March 2013): extending Glaisher (1900)},
  pdfauthor={Haobo Ma and Wenlin Zhang},
  pdfsubject={Exact literal binomial remainder counts and sharp prime-cube growth thresholds},
  pdfkeywords={A222759, A222760, binomial congruences, Glaisher, CC BY 4.0}
}
\newtheorem{theorem}{Theorem}
\newtheorem{lemma}[theorem]{Lemma}
\newtheorem{corollary}[theorem]{Corollary}
\newcommand{\pr}{\mathbb P}
\newcommand{\ind}{\mathbf 1}
\newcommand{\binrem}[3]{\left[\binom{#1#2}{#2}\right]_{#2^{#3}}}
\title{Exact counts and sharp growth for OEIS A222759, first comment
(12 March 2013): extending Glaisher (1900)}
\author{Haobo Ma\\\small ChronoAI Pte Ltd, Singapore
\and Wenlin Zhang\\\small National University of Singapore\\
\small ChronoAI Pte Ltd, Singapore}
\date{November 6, 2026}
\begin{document}
\maketitle
\begin{abstract}
Extending Glaisher's (1900) congruence to exact literal-remainder counts,
we determine OEIS A222759 and prove Noe's growth assertion in its first comment
(12 March 2013; revision 12, 2 November 2024).
For every positive integer $n$, its conjectured count equals
$\pi(\lfloor n^{1/3}\rfloor)$ plus two explicit corrections supported below $27$.
Consequently, for every $k\ge3$, the count is at least $k$ exactly when
$n\ge p_k^3$, sharpening the stated sufficient threshold $p_{p_k^3}$.
The same description determines A222760, including its exceptional relation
to A222759 at $n=5$. We give a self-contained proof, the counts modulo $p$
and $p^2$, and elementary consequences for the linked prime-counting entries.
\end{abstract}

\section{The count and its exact formula}
For a positive integer $m$, write $[x]_m$ for the least nonnegative remainder
of the integer $x$ modulo $m$. Let $p_j$ be the $j$th prime, with $p_1=2$,
and let $\pi(x)$ count the primes not exceeding the real number $x$.
Noe's entry A222759~\cite{noe}, first comment (12 March 2013), asks for
growth of a literal-remainder count; Theorem~\ref{thm:exact} determines it
and Corollary~\ref{cor:growth} gives its sharp threshold.
The primes counted form
\[
 W(n)=\left\{p\in\pr:\binrem{n}{p}{3}\ne n\right\},\qquad n\ge1.
\]
The comparison is with the integer $n$, not its residue class. The entry's
first comment says ``It appears that'' $a(n)\ge k$ for $k>2$ and
$n\ge p_{p_k^3}$. The cross-reference A096328~\cite{nested} confirms the
nested indexing. We first give a finite description of $W(n)$, so that
$a(n)=|W(n)|$ is well-defined.

Set
\begin{align*}
 E_2(n)&=\begin{cases}\{2\},&n<8\text{ and }n\equiv2,3\pmod4,\\
                         \varnothing,&\text{otherwise},\end{cases}\\
 E_3(n)&=\begin{cases}\{3\},&n<27\text{ and }n\equiv2\pmod3,\\
                         \varnothing,&\text{otherwise}.\end{cases}
\end{align*}
\begin{theorem}\label{thm:exact}
For every integer $n\ge1$, the following disjoint union is exact:
\begin{equation}\label{eq:set}
 W(n)=\{p\in\pr:p^3\le n\}\ \dot\cup\ E_2(n)\ \dot\cup\ E_3(n).
\end{equation}
In particular,
\begin{equation}\label{eq:count}
 a(n)=\pi(\lfloor n^{1/3}\rfloor)
 +\ind_{\{n<8,\ n\bmod4\in\{2,3\}\}}
 +\ind_{\{n<27,\ n\bmod3=2\}}.
\end{equation}
Thus $a(n)=\pi(\lfloor n^{1/3}\rfloor)$ for every $n\ge27$.
\end{theorem}

The classical input is Glaisher's 1900 congruence
$\binom{np}{p}\equiv n\pmod{p^3}$ for primes $p\ge5$.
Straub~\cite[\S2, p.~899]{straub} records this attribution and the harmonic
congruences underlying it; his equation (1) states the stronger
Ljunggren congruence $\binom{ap}{bp}\equiv\binom ab\pmod{p^3}$.
Here the additional step is to distinguish equality of literal remainders
from congruence, and to treat $p=2,3$ explicitly.

\section{A self-contained proof}
\begin{lemma}[Glaisher (1900), reproved]\label{lem:glaisher}
For every prime $p\ge5$ and integer $n\ge1$,
$\binom{np}{p}\equiv n\pmod{p^3}$.
\end{lemma}
\begin{proof}
All fractions below have denominators coprime to $p$; congruences of these
fractions mean congruences after inverting their denominators. Put
\[
 H_1=\sum_{j=1}^{p-1}\frac1j,\qquad
 H_2=\sum_{j=1}^{p-1}\frac1{j^2},\qquad
 S_2=\sum_{1\le i<j\le p-1}\frac1{ij}.
\]
Inversion permutes the nonzero residues modulo $p$, whence
\[
 H_2\equiv\sum_{j=1}^{p-1}j^2
 =\frac{p(p-1)(2p-1)}6\equiv0\pmod p.
\]
The denominator $6$ is invertible because $p\ge5$. Substituting $p-j$ for
$j$ and expanding modulo $p^2$ gives
\[
 2H_1=\sum_{j=1}^{p-1}\left(\frac1j+\frac1{p-j}\right)
 \equiv-pH_2\equiv0\pmod{p^2}.
\]
Hence $H_1\equiv0\pmod{p^2}$, and
$2S_2=H_1^2-H_2$ gives $S_2\equiv0\pmod p$.
Since $p$ is odd, the product formula for the binomial coefficient yields
\[
 \binom{np}{p}
 =n\prod_{j=1}^{p-1}\left(1-\frac{np}{j}\right)
 \equiv n\bigl(1-npH_1+n^2p^2S_2\bigr)
 \equiv n\pmod{p^3}.
\]
Every omitted product term contains at least three factors of $p$.
No division by $n$ is used, so the argument also covers $p\mid n$.
\end{proof}

\begin{proof}[Proof of Theorem~\ref{thm:exact}]
If $p^3\le n$, then $0\le\binrem{n}{p}{3}<p^3\le n$, so $p\in W(n)$,
including the boundary $n=p^3$.
If $p^3>n$ and $p\ge5$, Lemma~\ref{lem:glaisher} says that the remainder
is $n$, because $0\le n<p^3$. Such a prime is not in $W(n)$.

For $p=2$ the exact identity
\[
 \binom{2n}{2}-n=2n(n-1)
\]
is divisible by $8$ exactly when $4\mid n(n-1)$, or equivalently
$n\equiv0,1\pmod4$. When $n<8$, the other two residue classes give precisely
$E_2(n)$; when $n\ge8$, the range argument already applies.
For $p=3$,
\[
 \binom{3n}{3}-n=\frac92 n^2(n-1).
\]
Since $2$ is coprime to $3$, this is divisible by $27$ exactly when
$3\mid n^2(n-1)$, that is, $n\equiv0,1\pmod3$.
For $n<27$ the remaining class gives $E_3(n)$, and for $n\ge27$ the range
argument applies. The three sets in~\eqref{eq:set} are disjoint by their
defining inequalities, and all are finite. Counting proves~\eqref{eq:count}.
\end{proof}

\section{Growth and companion assertions}
\begin{corollary}[Sharp growth threshold]\label{cor:growth}
For integers $k\ge3$ and $n\ge1$,
\[
 a(n)\ge k\quad\Longleftrightarrow\quad n\ge p_k^3.
\]
In particular, Noe's sufficient condition $n\ge p_{p_k^3}$ holds.
\end{corollary}
\begin{proof}
For $n<27$, equation~\eqref{eq:set} gives $W(n)\subseteq\{2,3\}$, so
$a(n)\le2$. For $n\ge27$, equation~\eqref{eq:count} gives
$a(n)=\pi(\lfloor n^{1/3}\rfloor)$, which is at least $k$ exactly when
$p_k^3\le n$. Finally, $p_m\ge m+1$ for every $m\ge1$, since the $m$
distinct integers $p_1,\ldots,p_m$ all lie in $\{2,\ldots,p_m\}$.
Thus $p_{p_k^3}>p_k^3$.
\end{proof}

The sufficient direction has an even shorter proof independent of the
classical congruence: if $n\ge p_k^3$, the first $k$ primes satisfy
$\binrem{n}{p_j}{3}<p_j^3\le p_k^3\le n$. This proves the original growth
assertion without presupposing finiteness of the full witness set.
For example, the sharp thresholds for $k=3,4$ are $125,343$, whereas
A096328 gives $691,2309$ for the proposed sufficient thresholds.

The two remaining count levels can also be described completely:
\begin{equation}\label{eq:small-levels}
\begin{split}
 a(n)\ge1&\quad\Longleftrightarrow\quad n\notin\{1,4\},\\
 a(n)\ge2&\quad\Longleftrightarrow\quad n\ge27
 \text{ or }n\in\{2,8,11,14,17,20,23,26\}.
\end{split}
\end{equation}
These follow by substituting $1\le n<27$ in~\eqref{eq:count}; for
$n\ge27$, the primes $2,3$ both lie in $W(n)$. These levels were not
part of Noe's first-comment assertion.

\begin{corollary}[The eventual-good-prime sequence]\label{cor:companion}
The least prime $b(n)$ in A222760~\cite{companion} exists for every
$n\ge1$. With $M(n)=\max(W(n)\cup\{1\})$, it is the first prime greater
than $M(n)$. More explicitly,
\begin{equation}\label{eq:companion-all}
 b(n)=p_{a(n)+1}\quad(n\ne5),\qquad b(5)=5.
\end{equation}
For $n\ge27$ it is
\[
 p_{\pi(\lfloor n^{1/3}\rfloor)+1}.
\]
\end{corollary}
\begin{proof}
Equation~\eqref{eq:set} is a finite description of all bad primes. Every
prime above their maximum is good. Any smaller candidate prime has a bad
prime at or above it, and therefore fails the defining condition.
For $n\ge27$ the bad primes are exactly the primes with $p^3\le n$.
For $8\le n<27$, they are $\{2\}$ or $\{2,3\}$; for $1\le n<8$,
the only witness set which is not an initial segment of the primes is
$W(5)=\{3\}$. This proves~\eqref{eq:companion-all}.
\end{proof}

The case $n=2$ gives $W(2)=\{2,3\}$ and $b(2)=5$.
For $p\ge5$, its congruence is Wolstenholme's (1862), as recalled by
Straub~\cite[\S2]{straub}. The A222760 comment points to Weisstein's
MathWorld entry~\cite[eq.~(6)]{mathworld}, which records this consequence
and attributes its observation there to Noe in 2005. The general formula
above extends this recorded special case.

\begin{corollary}[The two smaller moduli]\label{cor:lower}
For $e=1,2$, the number of primes $p$ with $\binrem{n}{p}{e}\ne n$ is
respectively $\pi(n)$ and $\pi(\lfloor\sqrt n\rfloor)$.
These are the counts identified in A222759's second comment by its
cross-references to A000720~\cite{pi} and A056811~\cite{square}.
\end{corollary}
\begin{proof}
For an odd prime $p$, the same product gives
$\binom{np}{p}\equiv n\pmod p$. Also $H_1\equiv0\pmod p$ by pairing $j$
with $p-j$, so expansion through the linear term gives the congruence
modulo $p^2$ for every odd prime, including $3$.
For $p=2$, the identity $\binom{2n}{2}-n=2n(n-1)$ is divisible by $4$.
Thus for either exponent $e$, the literal remainder equals $n$ exactly
when $n<p^e$; if $n\ge p^e$, it is smaller than $n$. Counting these primes
proves the two formulas.
\end{proof}

\paragraph{Higher powers.}
The source's second comment (14 March 2013) also suggests a failure for
higher powers, without specifying an exact replacement count. A precise
obstruction is that the uniform congruence for primes $p\ge5$ cannot be
extended to any exponent
$e\ge4$: $n=2$, $p=5$ gives $\binom{10}{5}=252$, and
$[252]_{5^e}=252\ne2$ for every $e\ge4$.
The assertion of a uniform congruence at every higher exponent is
\emph{not asserted in the source; it fails for this example}.
This obstruction does not determine finiteness or the exact number of
exceptional primes for the larger exponents.

\section{Elementary consequences for the linked entries}
Let $c_e(n)=\pi(n^{1/e})$ for positive integers $e,n$. Thus $c_1$ is
A000720 and $c_2$ is A056811; the counts above are $c_1,c_2$ and,
for $n\ge27$, $c_3$.
\begin{corollary}\label{cor:linked}
The following descriptions hold.
\begin{enumerate}
\item $c_e(n)$ counts primes $p$ with $p^e\le n$, and prime factors
whose exponent in $\operatorname{lcm}(1,\ldots,n)$ is at least $e$.
For $e=2$ it also counts positive integers at most $n$ with exactly
three divisors, and is the maximum size of a pairwise coprime collection
of composite integers at most $n$.
\item For $e\ge1$, $c_e(n)=j$ exactly when
$p_j^e\le n<p_{j+1}^e$, with the $j=0$ interval $1\le n<2^e$.
As formal power series,
\[
 \sum_{n\ge1}c_e(n)x^n=\frac{\sum_{p\in\pr}x^{p^e}}{1-x}.
\]
In particular,
\[
 \sum_{n\ge1}a(n)x^n
 =\frac{\sum_{p\in\pr}x^{p^3}}{1-x}
   +x^2+x^3+x^6+x^7+x^2\sum_{j=0}^{8}x^{3j}.
\]
\item For fixed $e$, the prime number theorem gives
$c_e(n)\sim e n^{1/e}/\log n$ as $n\to\infty$. In particular,
$a(n)\sim3n^{1/3}/\log n$.
\end{enumerate}
\end{corollary}
\begin{proof}
The exponent of $p$ in the least common multiple is the largest $r$
with $p^r\le n$, proving the first two interpretations.
An integer has exactly three divisors precisely when it is the square
of a prime. Every composite $m\le n$ has a prime factor at most
$\sqrt m\le\sqrt n$; pairwise coprime composites have distinct such
factors. Their number is at most $c_2(n)$, attained by the prime squares.
The interval statement is the definition of prime counting.
For the generating function, sum $\sum_{n\ge p^e}x^n=x^{p^e}/(1-x)$
over primes; each coefficient involves finitely many primes.
The two correction indicators in~\eqref{eq:count} give the displayed
polynomial. Finally, substitute $t=n^{1/e}$ into
$\pi(t)\sim t/\log t$, recorded in A000720~\cite{pi}.
\end{proof}

For $e=2$, these interpretations and consequences reprove A056811's
comments and formulas: its original least-common-multiple and
three-divisor descriptions; Hurt's prime-square comment (2013);
Xie's pairwise-coprime maximum (2024); Herv\'e's constant blocks (2013);
Gutkovskiy's generating function (2019); and Richman's asymptotic (2022).
The same argument gives the partial sums of the prime indicator and
inverse-prime relations in A000720. The partitions of $2n$ into two
parts with smaller part prime are precisely $(p,2n-p)$ with $p\le n$,
giving Hurt's 2013 interpretation there. Also the primes dividing $n!$
are exactly those at most $n$, proving its factor-count description.

\section{Coverage and verification}
\begin{center}\small
\begin{tabular}{@{}>{\raggedright\arraybackslash}p{0.32\linewidth}>{\raggedright\arraybackslash}p{0.62\linewidth}@{}}
\toprule
Source item & Mathematical coverage\\
\midrule
A222759, first comment (2013) & Proved; Corollary~\ref{cor:growth} sharpens its sufficient threshold.\\
A222759, sharp converse and small levels & Not posed; proved in Corollary~\ref{cor:growth} and equation~\eqref{eq:small-levels}.\\
A222759, count definition & Determined by Theorem~\ref{thm:exact}.\\
Classical congruence & Reproved; prior proof Glaisher (1900), via Straub (2011), \S2.\\
A222759, second comment: $p,p^2$ & Reproved identifications; Noe (2013), with written proof in Corollary~\ref{cor:lower}.\\
A222759, higher-power remark & Obstruction proved for the uniform-congruence interpretation; full counting question not covered.\\
Uniform higher-power extension & Not asserted in the source; fails for $(n,p)=(2,5)$.\\
A222760, least-prime definition & Determined for every $n$; Corollary~\ref{cor:companion}.\\
A222760, relation to $a(n)$ & Not posed; proved in equation~\eqref{eq:companion-all}.\\
A222760, $n=2$ comment & Reproved; prior proof Wolstenholme (1862); Noe (2005) recorded in MathWorld, eq.~(6).\\
A096328 & Definition retained; gives a sufficient, not sharp, growth threshold.\\
A056811, elementary comments/formulas & Reproved interpretations, blocks, generating function and PNT consequence; Corollary~\ref{cor:linked}.\\
General $c_e$ and cubic generating function & Not posed; proved in Corollary~\ref{cor:linked}.\\
A000720, elementary identities & Proved consequences above, including its two-part partitions and factorial prime factors.\\
A000720, analytic bounds and remaining identities & Not covered; no implication here for RH or the Hardy--Littlewood conjecture.\\
\bottomrule
\end{tabular}
\end{center}
The declaration
\texttt{D5.S3.Arith.PrimeRemainders.}\allowbreak
\texttt{OeisA222759Growth.result}
states the original sufficient-threshold assertion: for $k\ge3$, $n\ge1$,
and $n\ge p_{p_k^3}$, there are $k$ distinct prime witnesses. Its
\href{https://github.com/the-omega-institute/trureturing/blob/3d781ae04219d5c99e87e82e3bbc3f6816e3691e/D5/S3/Arith/PrimeRemainders/OeisA222759Growth.lean}{pinned Lean source}
uses the first $k$ primes and the remainder-range argument.
The Lean statement has the original nested-prime hypothesis; it does not
state finiteness, an exact cardinality, or the sharp converse.
The exact count, its sharp threshold, the companion and smaller-modulus
formulas, and Corollary~\ref{cor:linked} have the written proofs above.

The accompanying Python program performs exact finite checks of all three
moduli for $1\le n\le512$ and all $25$ primes at most $97$, and of both
cubic-remainder sequences against their OEIS tables for $1\le n\le10000$.
These finite checks supplement the proofs for arbitrary positive $n$.
The program also checks the small count levels, companion relation,
generating-function coefficients and elementary linked interpretations
over the explicit ranges in \texttt{verification.json}.
The source archive contains \texttt{note.tex}, \texttt{verify.py},
\texttt{verification.json}, \texttt{metadata.json}, and rebuild instructions.

\begin{thebibliography}{9}
\bibitem{noe} T.~D. Noe, \emph{A222759}, The On-Line Encyclopedia of
Integer Sequences, March 12, 2013; revision 12, November 2, 2024.
\url{https://oeis.org/A222759}.
\bibitem{straub} A. Straub, A $q$-analog of Ljunggren's binomial congruence,
\emph{Discrete Mathematics \& Theoretical Computer Science Proceedings}
\textbf{AO} (2011), 897--902.
\href{https://doi.org/10.46298/dmtcs.2962}{doi:10.46298/dmtcs.2962};
arXiv:1103.3258v1.
\bibitem{companion} T.~D. Noe, \emph{A222760}, The On-Line Encyclopedia
of Integer Sequences, March 13, 2013; revision 9, February 16, 2025.
\url{https://oeis.org/A222760}.
\bibitem{nested} C. Hilliard, \emph{A096328}, The On-Line Encyclopedia
of Integer Sequences, August 2, 2004; revision 14, October 31, 2013.
\url{https://oeis.org/A096328}.
\bibitem{pi} The OEIS Foundation Inc., \emph{A000720: the prime-counting
function}, The On-Line Encyclopedia of Integer Sequences.
\url{https://oeis.org/A000720}.
\bibitem{square} Labos Elemer, \emph{A056811}, The On-Line Encyclopedia of
Integer Sequences, August 28, 2000; revision 49, August 15, 2024.
\url{https://oeis.org/A056811}.
\bibitem{mathworld} E.~W. Weisstein, \emph{Central Binomial Coefficient},
MathWorld, eq.~(6).
\url{https://mathworld.wolfram.com/CentralBinomialCoefficient.html}.
\end{thebibliography}
\medskip\noindent\small
\copyright\ 2026 Haobo Ma and Wenlin Zhang.
Licensed under \href{https://creativecommons.org/licenses/by/4.0/}{CC BY 4.0}.
\end{document}
